%%=============================================================================
%% CAPITULO 4 -- APRESENTACAO, ANALISE E DISCUSSAO DOS DADOS
%%=============================================================================

\section{Análise e discussão dos dados}

A análise dos dados buscará identificar padrões de evolução dos custos do SINAPI 
ao longo do tempo e diferenças entre localidades, avaliando de que forma essas 
informações podem contribuir para o planejamento de investimentos públicos. 
A exploração inicial permitirá caracterizar as séries históricas, 
observar variações territoriais e temporais e identificar aspectos relevantes para a 
posterior aplicação dos modelos estatísticos e de aprendizado de máquina. 

Os resultados dos modelos serão analisados de forma comparativa, considerando métricas quantitativas de 
desempenho e sua capacidade de apoiar decisões relacionadas a custos e investimentos públicos. 
A discussão buscará relacionar os resultados obtidos ao problema da pesquisa, 
especialmente quanto ao potencial da análise preditiva do SINAPI para antecipar 
variações de custos, apoiar o planejamento orçamentário e subsidiar políticas públicas 
baseadas em evidências.

\subsection{Apresentação dos resultados}

A análise será organizada nas dimensões apresentadas na 
Tabela~\ref{tab:dimensoes-analise-sinapi}, permitindo relacionar as características 
dos dados às diferentes etapas previstas para o desenvolvimento da pesquisa.

\begin{table}[htb] 
    \centering 
    \caption{Dimensões previstas para a análise dos dados do SINAPI} 
    \label{tab:dimensoes-analise-sinapi}
     \begin{tabular}{p{0.25\textwidth} p{0.60\textwidth}}
         \hline
          \textbf{Dimensão} & \textbf{Finalidade da análise} \\
           \hline
            Temporal & Identificar tendências e variações dos custos ao longo do tempo.\\
            Territorial & Comparar o comportamento dos custos entre diferentes localidades.\\
            Preditiva & Avaliar modelos estatísticos e de aprendizado de máquina para previsão de custos. \\
            Decisória & Examinar como os resultados podem apoiar o planejamento de investimentos públicos. \\
             \hline
             \end{tabular}
             \fonte{Elaboração própria.}
             \end{table}

\subsection{Discussão}

A discussão confrontará os resultados alcançados pelos diferentes modelos 
com o referencial teórico apresentado no Capítulo~2. Além da precisão preditiva, 
será considerada a possibilidade de utilização dos resultados como evidência para 
subsidiar decisões públicas, reconhecendo as limitações decorrentes da qualidade dos dados, 
do recorte temporal escolhido e das diferenças territoriais existentes na base SINAPI.
